\noindent
\textbf{CS4268 AY2025-26 Sem 2}
\\
\noindent
Lecturer: Divesh Aggarwal
\hfill
Lecture 8                      %%% FILL IN LECTURE NUMBER HERE
\hfill
Date: March 12, 2026          %%% FILL IN LECTURE DATE HERE


\noindent
\rule{\textwidth}{1pt}

\medskip


\section{Simon's Problem and Algorithm}

\subsection[\texorpdfstring{Simon over $\bm{\mathbb{Z}_2^n}$}{}]{Simon over $\bm{\mathbb{Z}_2^n}$} 

\paragraph{Problem:} Given $f: \{0,1\}^n \rightarrow \{0,1\}^n$ where it is promised that $\exists s \in \mathbb{Z}_2^n$ such that $\forall x,y$: $f(x)=f(y) \iff x=y\oplus s$, find $s$.

One way to view this is as a \emph{hidden period problem} over the group $\mathbb{Z}_2^n$. The usual notion of periodicity is $f(x)=f(x+r)$ for some hidden shift $r$ under ordinary addition; here the ``addition'' is bitwise XOR, so the period is the hidden string $s$ satisfying $f(x)=f(x \oplus s)$. Algebraically, the hidden structure is captured by the subgroup $\{0^n, s\} \leq \mathbb{Z}_2^n$: two inputs are equivalent (i.e.\ $f$-indistinguishable) if and only if they lie in the same coset of this subgroup. When $s=0^n$ the function is one-to-one; when $s\neq 0^n$ every output has exactly two preimages, namely $x$ and $x \oplus s$.

\paragraph{Classical Protocols}
\begin{enumerate}
    \item A randomized algorithm: randomly pick $k$ (TBD) distinct inputs $x_i$ and query to get $f(x_i)$. If there is no collision, answer $s=0^n$. If there is a collision $f(x_i)=f(x_j)$, then answer $s=x_i \oplus x_j$. 
    
    \item Analysis: if actually $s=0^n$ then the algorithm works with certainty. If actually $s \neq 0^n$ then we might answer wrongly if there are no observed collisions. Let the probability that the samples have no collision after $k\geq 2$ queries be $P(k)$. Suppose after $j-1$ queries there is no matching. The probability that on the $j$th, i.e. next query, there is still no collision with one of the previous queries is then $1-\frac{j-1}{2^n-(j-1)}$. Thus after $k$ rounds,
    \begin{align*}
        P(k) &= \prod_{j=2}^k \left[ 1- \frac{j-1}{2^n-(j-1)} \right]\\
        &\geq 1-\sum_{j=2}^k \left[ \frac{j-1}{2^n-(k-1)} \right]\\
        &\geq 1- \frac{\Theta(k^2)}{2^n-\Theta(k)}.
    \end{align*}
    Remember that we \textit{do} want a collision! So we want $P(k)$ to be small, say $\varepsilon$ (equivalently, coll. prob. is $1-\varepsilon$). Then $P(k)<\varepsilon \implies k \in \Omega(2^{n/2})$. Here in the first $\geq$ we used $(1-a)(1-b)\geq 1-(a+b)$ for $0\leq a,b \leq 1$. Also note that we can assume from the outset that $k\leq 2^{n/2}$ because if $k>2^{n/2}$ then there has to be a collision.
    
    \item Simon also showed \textit{any} randomized algorithm needs $k=\Omega(2^{n/2})$ queries (proof is similar to the one above).
\end{enumerate}
    
\paragraph{Simon's Quantum Algorithm}
\begin{enumerate}
    \item Quantum subroutine: Start with $n$ work qubits and $n$ ancilla qubits $\ket{0^n}\ket{0^n}$. Apply $$(H^{\otimes n}\otimes I)O_f(H^{\otimes n} \otimes I)$$ then measure the work qubits (first register).

    \begin{figure}[!ht]
    \centering
    \begin{quantikz}
    \lstick{$\ket{0^n}$} & \gate{H^{\otimes n}} & \gate[wires=2]{O_f} & \gate{H^{\otimes n}} & \meter{} \\
    \lstick{$\ket{0^n}$} & \qw & \qw & \qw & \qw
    \end{quantikz}
    \caption{Simon's quantum subroutine for $\mathbb{Z}_2^n$.}
    \label{fig:simon_Z2n}
    \end{figure}

    To see what this circuit does, it helps to track the state at each stage:
    \begin{itemize}
        \item \textbf{After first $H^{\otimes n}$:} The first register is in the uniform superposition over all $n$-bit strings,
        \[
            \frac{1}{\sqrt{2^n}} \sum_{x \in \{0,1\}^n} \ket{x}\ket{0^n}.
        \]
        \item \textbf{After oracle $O_f$:} The oracle entangles the two registers, writing $f(x)$ into the second register for each $x$ simultaneously:
        \[
            \frac{1}{\sqrt{2^n}} \sum_{x \in \{0,1\}^n} \ket{x}\ket{f(x)}.
        \]
        \item \textbf{After measuring the second register:} Suppose the outcome is some value $z = f(x_0)$. Because $f$ is two-to-one (for $s \neq 0^n$), exactly two inputs $x_0$ and $x_0 \oplus s$ map to $z$. The first register collapses to the \emph{pair state}
        \[
            \frac{1}{\sqrt{2}}\bigl(\ket{x_0} + \ket{x_0 \oplus s}\bigr).
        \]
        This state encodes the hidden string $s$: the two branches differ by exactly $s$ under XOR. Note that $x_0$ is random and unknown, but that does not matter for the next step.
        \item \textbf{After second $H^{\otimes n}$:} Applying the Hadamard transform to the pair state produces, for each $y \in \{0,1\}^n$, an amplitude proportional to
        \[
            (-1)^{x_0 \cdot y}\bigl(1 + (-1)^{s \cdot y}\bigr),
        \]
        where $x \cdot y = \sum_i x_i y_i \pmod{2}$ is the bitwise dot product. The bracket $1+(-1)^{s \cdot y}$ is either $2$ (when $s \cdot y = 0$) or $0$ (when $s \cdot y = 1$). Thus only those $y$ with $s \cdot y = 0 \pmod{2}$ have nonzero amplitude, and measurement of the first register must yield such a $y$.
    \end{itemize}

    \item Classical postprocessing: repeat the quantum subroutine $n-1$ times to get $n-1$ strings $y_i$ all satisfying $y_i \cdot s =0$, where the $y_i$'s are linearly independent with positive probability. If indeed they are all linearly independent, then we use Gaussian elimination to solve the system, obtaining two solution candidates $0^n$ and $s' \neq 0^n$. 
    
    Compare the values of $f(0^n)$ and $f(s')$. If $f(0^n)=f(s')$, answer $s=s'$; otherwise $f(0^n)\neq f(s')$, answer $s=0^n$.
    
    \item Boosting: repeat the entire process $k$ times to boost the probability of obtaining linearly independent $y_i$ (linear independence can be easily checked). If for \text{any} run $f(0^n)=f(s')$, answer $s=s'$; otherwise if for \text{all} runs $f(0^n)\neq f(s')$, answer $s=0^n$.
\end{enumerate}
    
\paragraph{Analysis of Simon's algorithm}
\begin{enumerate}
    \item Quantum: after applying the quantum circuit (before measurement) the state is (good exercise to verify yourselves)
    \begin{align*}
        \ket{\psi} = \sum_y \ket{y} \left( \frac{1}{2^n} \sum_x (-1)^{x \cdot y} \ket{f(x)} \right).
    \end{align*}
    To see where the measurement constraint $y \cdot s = 0$ comes from, it is instructive to track the amplitude of a fixed $\ket{y}$ through the final Hadamard. After the second register is measured and found to equal $z = f(x_0)$, the first register is in the pair state $\frac{1}{\sqrt{2}}(\ket{x_0}+\ket{x_0 \oplus s})$. The Hadamard identity $H^{\otimes n}\ket{x} = \frac{1}{\sqrt{2^n}}\sum_{y}(-1)^{x \cdot y}\ket{y}$ gives
    \begin{align*}
        H^{\otimes n} \cdot \frac{1}{\sqrt{2}}\bigl(\ket{x_0}+\ket{x_0 \oplus s}\bigr)
        &= \frac{1}{\sqrt{2 \cdot 2^n}} \sum_y (-1)^{x_0 \cdot y}\bigl(1+(-1)^{s \cdot y}\bigr)\ket{y}.
    \end{align*}
    The factor $1+(-1)^{s \cdot y}$ is $2$ if $s \cdot y = 0 \pmod{2}$ and $0$ if $s \cdot y = 1 \pmod{2}$: this is \emph{destructive interference}, where the two branches of the superposition cancel whenever $s \cdot y \neq 0$. Consequently only those $y$ satisfying $s \cdot y = 0$ survive measurement.

    If $s=0^n$, $f$ is injective so $p(y) = \frac{1}{2^n}$ for all $y$.

    If $s \neq 0^n$, $f$ is two-to-one, so if $z \in \text{range}(f)$ then $f(x_z) = f(x_z') = z$ for $x_z = x_z' \oplus s$. Therefore in this case
    \begin{align*}
    p(y)=
    \begin{cases}
        \frac{1}{2^{n-1}} & \text{if } y\cdot s = 0\\
        0 & \text{if } y\cdot s = 1.
    \end{cases}
    \end{align*}

    \item The transform $H^{\otimes n}$ appearing here is the \emph{Walsh--Hadamard Transform (WHT)}, which is the Fourier transform for the group $(\mathbb{Z}_2)^n$. Concretely, the \emph{Walsh characters} (i.e.\ the irreducible unitary characters of $(\mathbb{Z}_2)^n$) are the functions
    \[
        \varepsilon_s : x \mapsto (-1)^{s \cdot x}, \quad s \in \{0,1\}^n,
    \]
    which form a complete orthogonal basis for functions $\{0,1\}^n \to \mathbb{R}$ (or $\mathbb{C}$). The WHT decomposes any such function into its Walsh-character components, exactly as the classical Fourier transform decomposes a function into complex exponentials. The formula
    \[
        H^{\otimes n}\ket{x} = \frac{1}{\sqrt{2^n}}\sum_{y}(-1)^{x \cdot y}\ket{y}
    \]
    is precisely the change of basis from the computational basis to the Walsh-character basis.

    From this viewpoint, the interference argument above has a clean spectral interpretation. Measuring the second register collapses the first register to the \emph{indicator function} of the two-element set $\{x_0, x_0 \oplus s\}$. The WHT of this function is non-zero only on the subspace
    \[
        s^\perp = \{y \in \{0,1\}^n : y \cdot s = 0 \pmod{2}\},
    \]
    which has dimension $n-1$ (i.e.\ $2^{n-1}$ elements). In other words, the WHT concentrates all probability mass on $s^\perp$, and the measurement outcomes are the ``Fourier frequencies'' of the hidden period subgroup $\{0^n, s\}$. Collecting $n-1$ linearly independent elements of $s^\perp$ and solving the resulting linear system over $\mathbb{F}_2$ recovers $s$ — this is the classical Gaussian elimination step.

    \item Postprocessing: suppose we have linearly independent $\{y_i\}_{i=1}^{n-1}$. Gaussian elimination gives the two possible solutions: $0^n,s'\neq 0^n$. Suppose actually $s=0^n$. Then $f(0^n)\neq f(s')$, so our answer $s=0^n$ is always correct. Suppose actually $s \neq 0^n$. Then $f(s)=f(0^n)=f(s')$, so our answer $s'=s$ is also always correct.
    
    \item Without the linear independence of $\{y_i\}_{i=1}^{n-1}$ the postprocessing fails, thus we need to ensure the probability of linear independence is nonzero (and preferably high).
    
    Probability of linear independence: given $m$ LI strings $y_1,\dots,y_m$, spanning subspace is size $2^m$. Prob of $y_{m+1}$ being LI to this is $\frac{2^n-2^m}{2^n}$. So for a single run, the probability of $n-1$ LI strings is $P=(1-\frac{2^1}{2^n})\dots(1-\frac{2^{n-2}}{2^n}) \geq 1-\sum_{i=1}^{n-2} \frac{1}{2^i} > \frac{1}{2}$. 
    
    Thus, the probability of \textit{no} LI after $k$ runs is at most $(1-\frac{1}{2})^k < e^{-k/2}$, or equivalently, probability of \textit{at least one} LI after $k$ runs is at least $1 - e^{-k/2}$. This is really just the classical boosting via majority vote trick.
\end{enumerate}



\paragraph{Worked Example ($n=3$, $s=101$)}

Suppose $n=3$ and the hidden string is $s=101$. Each run of the quantum subroutine outputs a random $y=(y_1,y_2,y_3)$ satisfying
\[
    y \cdot s = y_1 \cdot 1 + y_2 \cdot 0 + y_3 \cdot 1 = y_1 + y_3 = 0 \pmod{2}.
\]
The valid outcomes are exactly those strings with $y_1 = y_3$, namely
\[
    \{000,\; 010,\; 101,\; 111\}.
\]
After collecting two linearly independent equations, say $y^{(1)} = 010$ (giving $0=0$, trivially satisfied, so this is not useful independently) and $y^{(2)} = 101$ (giving $1+1=0 \pmod 2$, consistent with $s=101$), Gaussian elimination over $\mathbb{F}_2$ yields $s=101$ as the unique nonzero solution. One then verifies by querying $f(0^3)$ and $f(101)$ and confirming they match.

\paragraph{Remark: Connection to Shor's Algorithm}

Simon's algorithm is historically significant as the first example of a \emph{super-polynomial} quantum speedup over any classical algorithm for a promise problem. More importantly, it establishes a design pattern that recurs in Shor's algorithm:
\begin{enumerate}
    \item Create a uniform superposition over all inputs.
    \item Query a function with hidden periodic structure.
    \item Apply a Fourier-type transform (here $H^{\otimes n}$, i.e.\ the WHT; in Shor, the $\mathrm{DFT}_N$).
    \item Use interference to concentrate measurement outcomes on constraints satisfied by the hidden period.
    \item Recover the period via classical postprocessing (here Gaussian elimination; in Shor, the continued-fractions algorithm).
\end{enumerate}
In both cases the quantum computer does not output the period directly --- it outputs \emph{constraints} on the period, and the classical step reconstructs it. Simon over $\mathbb{Z}_N$ (see below) makes this analogy even more explicit.

The same 5-step skeleton appears across several quantum algorithms, each one being an instance of the \emph{hidden subgroup problem} (HSP) for a different group:

\begin{center}
\begin{tabular}{llll}
\hline
\textbf{Algorithm} & \textbf{Group} & \textbf{Hidden structure} & \textbf{Fourier transform} \\
\hline
Bernstein--Vazirani & $(\mathbb{Z}_2)^n$ & Hidden linear function ($f(x)=s\cdot x$) & WHT ($H^{\otimes n}$) \\
Simon's            & $(\mathbb{Z}_2)^n$ & Hidden subgroup $\{0^n,s\}$ (XOR period)  & WHT ($H^{\otimes n}$) \\
Simon over $\mathbb{Z}_N$ & $\mathbb{Z}_N$ & Hidden subgroup $L\mathbb{Z}$ (additive period) & $\mathrm{DFT}_N$ \\
Shor's             & $\mathbb{Z}_{2^n}$ & Multiplicative order $r = \mathrm{ord}_N(a)$ & $\mathrm{QFT}_{2^n}$ \\
\hline
\end{tabular}
\end{center}

\noindent The key distinction is the underlying group: Bernstein--Vazirani and Simon operate over the binary group $(\mathbb{Z}_2)^n$ under XOR, so the correct Fourier transform is the WHT (whose characters are the Walsh functions $\varepsilon_s(x) = (-1)^{s \cdot x}$). Shor operates over $\mathbb{Z}_{2^n}$ under ordinary addition, so the correct transform is the $\mathrm{QFT}$ (whose characters are $\varepsilon_s(x) = e^{2\pi isx/N}$). The algorithms look structurally identical because they are: the only difference is which Fourier transform is the right one for the group at hand.

\subsection[\texorpdfstring{Simon over $\bm{\mathbb{Z}_N}$}{}]{Simon over $\bm{\mathbb{Z}_N}$} 

\paragraph{Problem:} Given $f: \mathbb{Z}_N \rightarrow \mathbb{Z}_N$ where it is promised that $\exists L \in \mathbb{Z}_N$ such that $\forall x,y$: $f(x)=f(y) \iff x=y+L$, find $L$. Here $N=2^n$ is a power of two.

For simplicity we consider only the case where $L$ divides $N$ (ruling out $L=0$ in particular).

\paragraph{Simon's Quantum Algorithm}
\begin{enumerate}
    \item Quantum subroutine: Start with $n$ work qubits and $n$ ancilla qubits $\ket{0^n}\ket{0^n}$. Apply $$(\mathsf{DFT_N}\otimes I)O_f(\mathsf{DFT_N} \otimes I)$$ then measure the work qubits (first register).

    \begin{figure}[!ht]
    \centering
    \begin{quantikz}
    \lstick{$\ket{0^n}$} & \gate{\mathsf{DFT_N}} & \gate[wires=2]{O_f} & \gate{\mathsf{DFT_N}} & \meter{} \\
    \lstick{$\ket{0^n}$} & \qw & \qw & \qw & \qw
    \end{quantikz}
    \caption{Simon's quantum subroutine for $\mathbb{Z}_N$.}
    \label{fig:simon_ZN}
    \end{figure}

    \item Classical postprocessing: for two runs of the quantum subroutine above we obtain the values $iN/L,jN/L$. Use this to compute our estimate $\hat{L} = N/\gcd(\frac{iN}{L},\frac{jN}{L})$.

    \item Boosting: run the quantum subroutine $k$ times (instead of just 2). From the $k$ measurement results $\{\frac{i_\alpha N}{L}\}_{\alpha=1}^k$ we obtain $\gcd\left(\frac{i_1N}{L},\dots,\frac{i_kN}{L}\right)$. Our estimate is correspondingly $\hat{L} = N/\gcd\left(\frac{i_1N}{L},\dots,\frac{i_kN}{L}\right)$.
    
\end{enumerate}
    
\paragraph{Analysis of Simon's algorithm}
\begin{enumerate}
    \item Quantum: The protocol is exactly the same as that for $\mathbb{Z}_2^n$, except that instead of $H^{\otimes n}$ we now use the $\mathsf{DFT_N}$, which is (before measurement) we get the state (good exercise to verify yourselves)
    \begin{align*}
        \ket{\psi} = \sum_{y=0}^{N-1} \ket{y} \left( \frac{1}{N} \sum_{x=0}^{N-1} e^{2\pi i xy/N} \ket{f(x)} \right).
    \end{align*}
    Note the analogy between this state and the post-circuit state in the $\mathbb{Z}_2^n$ case above. We now proceed to simplify this expression, using the periodicity property of $f$: $f(x)=f(x+tL)$ for $t=0,1,\dots,N/L-1$. We have
    \begin{align*}
        \ket{\psi} &= \sum_{y=0}^{N-1} \ket{y} \left( \frac{1}{N} \sum_{x=0}^{N-1} e^{2\pi i xy/N} \ket{f(x)} \right)\\
        &= \sum_{y=0}^{N-1} \ket{y} \left( \frac{1}{N} \sum_{x=0}^{L-1} \sum_{t=0}^{N/L-1} e^{2\pi i y(x+tL)/N} \ket{f(x+tL)} \right)\\
        &= \sum_{y=0}^{N-1} \ket{y} \left( \frac{1}{N} \sum_{x=0}^{L-1} e^{2\pi i xy/N} \ket{f(x)} \sum_{t=0}^{N/L-1} e^{2\pi i ytL/N} \right)\\
        &= \sum_{y=0,\frac{N}{L},\dots,(L-1)\frac{N}{L}} \ket{y} \left( \frac{1}{L} \sum_{x=0}^{L-1} e^{2\pi i xy/N} \ket{f(x)} \right)
    \end{align*}
    In the second equality we have simply regrouped terms; in the third equality we have invoked periodicity; in the last equality we have evaluated the geometric series
    \begin{align*}
        \sum_{t=0}^{N/L-1} e^{2\pi i ytL/N}=
        \begin{cases}
            \frac{N}{L} & \text{if } y=0 \text{ mod } \frac{N}{L}\\
            0 & \text{otherwise.}
        \end{cases}
    \end{align*}
    Measuring the work (first) register now gives $p(y)=\frac{1}{L}$ for those $y$ satisfying $y=0$ mod $\frac{N}{L}$, or equivalently $y=0,\frac{N}{L},\frac{2N}{L},\dots, \frac{(L-1)N}{L}$. That is,
    \begin{align*}
    p(y)=
    \begin{cases}
        \frac{1}{L} & \text{if } y = 0 \text{ mod } \frac{N}{L}\\
        0 & \text{otherwise.}
    \end{cases}
    \end{align*}
    As in the $\mathbb{Z}_2^n$ case above, what Simon really does is to refocus the probability mass onto the states $\ket{0},\ket{N/L},\ket{2N/L},\dots,\ket{(L-1)N/L}$, so that the final measurement yields a multiple of $N/L$. That is, a single run of this quantum circuit yields a number of the form $iN/L$, for some $i \in \{0,1,\dots,L-1\}$. Please note that we do not obtain the value $i$ -- rather it is the `overall' value $iN/L$ that we receive. From this we shall attempt to extract $L$. 
    
    \item Postprocessing: our estimate for the period $L$ is $\hat{L} = N/\gcd(\frac{iN}{L},\frac{jN}{L})$. Since $\gcd(i,j)=1 \iff \gcd(\frac{iN}{L},\frac{jN}{L})= N/L$, if it is indeed the case that $\gcd(i,j)=1$, then we get $\hat{L}=L$, yielding the correct answer.
    
    \item The classical postprocessing fails if $\gcd(i,j) > 1$. Thus we need to ensure the probability that $\gcd(i,j)=1$ is nonzero. To do this we turn to some number theory. For any prime $p$ we have $\Pr[p|i] = \lfloor N/p \rfloor/N \leq 1/p$. By independence we have $\Pr[p|i \text{ and } p|j] = \Pr[p|i] \cdot \Pr[p|j] \leq \frac{1}{p^2}$. Since $\gcd(i,j) > 1$ implies that \textit{at least} a prime $p$ divides both $i,j$, we have by the union bound
    $$\Pr[\gcd(i,j) > 1] \leq \sum_{p \text{ prime}} \frac{1}{p^2} = \frac{1}{2^2} + \frac{1}{3^2} + \frac{1}{5^2} + \cdots.$$
    We can bound this sum as follows:
    \begin{align*}
    \sum_{p \text{ prime}} \frac{1}{p^2}
    &\leq \frac{1}{2^2} + \frac{1}{3^2} + \frac{1}{4^2} + \frac{1}{5^2} + \cdots\\
    &= \frac{1}{4} + \left(\frac{1}{3^2} + \frac{1}{4^2} + \frac{1}{5^2}\right) + \left(\frac{1}{6^2} + \cdots + \frac{1}{11^2}\right) + \left(\frac{1}{12^2} + \cdots + \frac{1}{23^2}\right) + \cdots\\
    &\leq \frac{1}{4} + \underbrace{\left(\frac{1}{3^2} + \frac{1}{3^2} + \frac{1}{3^2}\right)}_{3\text{ terms}} + \underbrace{\left(\frac{1}{6^2} + \cdots + \frac{1}{6^2}\right)}_{6\text{ terms}} + \underbrace{\left(\frac{1}{12^2} + \cdots + \frac{1}{12^2}\right)}_{12\text{ terms}} + \cdots\\
    &= \frac{1}{4} + \frac{3}{3^2} + \frac{6}{6^2} + \frac{12}{12^2} + \cdots\\
    &= \frac{1}{4} + \frac{1}{3}\left(1 + \frac{1}{2} + \frac{1}{4} + \cdots\right)\\
    &= \frac{1}{4} + \frac{1}{3}\cdot 2 = \frac{11}{12}.
    \end{align*}
    Therefore $\Pr[\gcd(i,j)=1] \geq 1/12$. $\frac{1}{12}$ is not great but this is no problem because simply by running the subroutine $k$ times we exponentially increase the success probability. This follows from the exact same logic above for the two-run-scenario. For $k$ runs, we have
    \begin{align*}
        \Pr \left[\gcd\left(i_1,\dots,i_k \right) >1 \right] \leq \sum_{p \text{ prime}} \frac{1}{p^k} = O(2^{-k}),
    \end{align*}
    since the dominant contribution comes from the prime $p=2$. Equivalently,
    \begin{align*}
        \Pr \left[\gcd\left(i_1,\dots,i_k \right) = 1 \right] \geq 1- \sum_{p \text{ prime}} \frac{1}{p^k} = 1- O(2^{-k}).
    \end{align*}

\end{enumerate}

